\section{Problem definition}
\label{sec:problem}

Let \(I\) be an IFC model, \(K\) a set of construction-domain rules, \(R\) a
robot description, and \(T\) a task specification. The compiler produces an
OpenUSD scene \(S\) and a provenance ledger \(P\):
\begin{equation}
  C(I,K,R,T) \rightarrow (S,P).
\end{equation}

The scene is composed from three fidelity layers:
\begin{equation}
  S = S_{\mathrm{global}} \oplus S_{\mathrm{workzone}}
  \oplus S_{\mathrm{interaction}},
\end{equation}
where \(\oplus\) denotes OpenUSD composition. The active task determines which
work-zone and interaction payloads are loaded.

A construction scene contract is a versioned set of requirements
\(\mathcal{Q}=\{q_1,\ldots,q_n\}\). A deterministic validator returns:
\begin{equation}
  V(S,T,\mathcal{Q}) \rightarrow E,
\end{equation}
where \(E\) is a structured set of violations containing a rule identifier,
severity, scene path, evidence, and allowed repair tools.

For every critical generated property \(p_i\), the compiler stores:
\begin{equation}
  p_i = (v_i, s_i, c_i, u_i),
\end{equation}
where \(v_i\) is the value, \(s_i\) its source, \(c_i\in[0,1]\) a confidence
score, and \(u_i\) an uncertainty description or interval.

Repair is then a constrained selection, not a probability maximization, with
two distinct acceptance levels. Each violation \(e\in E\) carries a small
allowed-tool set \(\mathcal{A}(e)\subseteq\mathcal{A}\) registered by its
rule. Writing \(E_0\), \(E_1\) for the violation sets before and after a
candidate fix \(a\in\mathcal{A}(e)\), and \(B(\cdot)\) for the blocking
(critical) members, a single repair step is \emph{submitted} only if it
resolves its target and adds no new blocking violation, so a no-op or a
mutation that merely hides the original violation cannot be accepted:
\begin{align}
  \mathrm{submit}(a,e) \iff {}& e\notin E_1\ \land\nonumber\\
  & \bigl[B(E_1)\setminus B(E_0)=\varnothing\bigr]\ \land\nonumber\\
  & \mathrm{EvidenceOK}(a,P) .
\end{align}
This is intentionally local: in a compound-fault scene, fixing \(e_1\) while
an unrelated \(e_2\) is still open does not block the step. The repaired
scene is then \emph{accepted} only when no blocking violation remains:
\begin{equation}
  \mathrm{accept}(S') \iff B\!\left(E_{\mathrm{final}}\right)=\varnothing,
\end{equation}
where a violation's identity is its rule and entity, not its diagnostic text.
Unsupported or safety-critical mutations are escalated to human review. The
whitelisted tool registry and this two-level re-validation gate are the
implemented mechanism; no probability is estimated and no such optimization
is solved at runtime.



\Cref{tab:info-mapping} organizes the contract families by the robot function
they serve, the task phase in which they act, and the measured consequence of
their absence or error. Three levels are distinguished throughout: whether a
building-information item can be \emph{expressed} in the scene, whether the
executor actually \emph{uses} it, and whether its corruption measurably
\emph{affects} the task outcome.

% 建筑信息 → 机器人功能 → 任务阶段 → 后果 映射表（插入 03-problem-definition 或 06）
% 注意: "实验证据"列指向具体结果小节; "未验证"如实标注

\paragraph{Provenance of the task information.}
\Cref{tab:info-provenance} states, for each field the executor consumes, where
the value comes from and who is responsible for it. Four source classes are
distinguished: \emph{IFC-native} values already present in the delivery,
\emph{derived} values computed from IFC geometry or material, \emph{task
annotation} supplied with the robot--task specification (the CSR property-set
layer), and \emph{assumed} defaults that must be escalated rather than trusted
silently. The mapping makes explicit that the pipeline does not pretend that
every required quantity is IFC-native: interface coordinates are extracted
from the IFC hole feature and then certified against an evidence record;
clearance and state are annotation responsibilities; mass is derived with an
explicit uncertainty; and an unannotated field is by construction assigned to
the escalation path, not to a guessed default (\Cref{sec:results}).

